\documentclass[12pt]{article}
\usepackage[margin=1in]{geometry}
\usepackage{xcolor}

\newcommand{\revised}[1]{\textcolor{blue}{#1}}

\begin{document}

\noindent Dear Editor,

\medskip

We thank the reviewer for their careful and constructive report. We have revised the manuscript to address all 5 major and 8 minor issues raised. All changes are text revisions; no new experiments were required. Below we provide point-by-point responses.

\bigskip

\section*{Response to Referee}

\subsection*{Major Issue 1 (REF-001): Rare-class narrative overclaimed}

\textit{``The introduction states `yet rare-class performance declines' without qualification. The conclusion converts an underpowered observation into a causal narrative.''}

\medskip

We agree completely. The rare-class comparison has aggregate power of only 0.20, and the data provide insufficient evidence to determine even the direction of the effect.

\textbf{Changes:} We have harmonized all rare-class language throughout the manuscript. The introduction now reads ``yet the rare-class comparison is statistically underpowered (aggregate power $= 0.20$), yielding insufficient evidence to determine whether ViT improves or degrades rare-class performance.'' The conclusion now frames the result as ``the rare-class comparison remains statistically underpowered, leaving insufficient evidence to determine whether CNN's inductive bias provides a meaningful advantage.'' The Discussion (Sec~4.1) conditions the inductive-bias interpretation on ``future multi-seed studies.''

\subsection*{Major Issue 2 (REF-002): Novelty claim too broad}

\textit{``Wu et al.\ (CQG 42, 165015) report class-wise accuracy tables. The novelty claim should be narrowed.''}

\medskip

We agree. Wu et al.\ report per-class metrics but with a different architecture and without paired statistical testing.

\textbf{Changes:} We have narrowed the claim to ``identical-recipe paired comparison with statistical testing of per-class differences and simulation-based power analysis'' and added a clarifying sentence acknowledging Wu et al.'s per-class results (Introduction).

\subsection*{Major Issue 3 (REF-003): Missing canonical review citation}

\textit{``Cuoco et al.\ (2025) Living Reviews in Relativity is missing. Wu et al.\ should be updated to the published CQG version.''}

\medskip

\textbf{Changes:} We have added Cuoco et al.\ (2025) to the introduction and updated Wu et al.\ to CQG 42, 165015.

\subsection*{Major Issue 4 (REF-004): Asymmetric hedging of interpretations}

\textit{``The Power\_Line self-attention interpretation is presented as explanatory while the rare-class interpretation is hedged as speculative. Neither has direct evidence.''}

\medskip

We agree. We have applied consistent speculative qualification throughout.

\textbf{Changes:} The Power\_Line interpretation now reads ``We hypothesize that this benefits classes with spatially distributed, multi-frequency structure\ldots but we lack direct evidence (attention maps or ablation studies) to confirm this mechanism.'' We also cite Raghu et al.\ (2021) for context on ViT representation structure.

\subsection*{Major Issue 5 (REF-005): CW abstract conflates veto efficiency with duty cycle}

\textit{``The abstract reports `approximately equal' veto efficiency without mentioning the CNN duty-cycle advantage.''}

\medskip

\textbf{Changes:} The abstract now includes the duty-cycle difference: ``though CNN retains a statistically significant duty-cycle advantage ($\Delta_\mathrm{DC} = -0.051$).'' We have also replaced ``first quantitative assessment'' with ``first proxy-level comparison'' in the introduction.

\subsection*{Minor Issues (REF-006 through REF-013)}

\begin{enumerate}
  \item[6.] \textbf{REF-006:} Corrected ViT O4 macro-F1 from 0.670 to 0.669 (double-rounding fix).
  \item[7.] \textbf{REF-007:} Changed ``all 8 rare classes'' to ``all 8 classes with small test sets.''
  \item[8.] \textbf{REF-008:} Softened Sec~3.3 heading to ``No Significant Evidence That Architecture Preference Is a Monotonic Function of Sample Size.''
  \item[9.] \textbf{REF-009:} Removed ad-hoc 20\% pass/fail threshold; now reports asymmetric degradation directly.
  \item[10.] \textbf{REF-010:} Added He \& Garcia (2009) and Buda et al.\ (2018) citations; reframed forbidden-proxy discussion as GW demonstration of known principle.
  \item[11.] \textbf{REF-011:} Added Raghu et al.\ (2021) citation for ViT interpretability context.
  \item[12.] \textbf{REF-012:} Added footnote to Table~I explaining one-sided test direction and two-sided $p$-value.
  \item[13.] \textbf{REF-013:} Updated Wu et al.\ citation from arXiv to CQG 42, 165015.
\end{enumerate}

\section*{Summary of Changes}

\subsection*{Major changes}
\begin{enumerate}
  \item Harmonized rare-class language throughout: all directional claims replaced with ``insufficient evidence'' framing (Abstract, Introduction, Discussion, Conclusion).
  \item Narrowed novelty claim to specify identical-recipe paired comparison with statistical testing and power analysis.
  \item Added Cuoco et al.\ (2025) Living Reviews in Relativity and updated Wu et al.\ to published CQG version.
  \item Applied consistent speculative qualification to all inductive-bias interpretations.
  \item Abstract now reports both veto efficiency and duty-cycle difference for CW analysis.
\end{enumerate}

\subsection*{Minor changes}
\begin{enumerate}
  \item Corrected double-rounding error (ViT O4 macro-F1: 0.669).
  \item Fixed ``8 rare classes'' labeling inconsistency.
  \item Softened Sec~3.3 heading.
  \item Removed unmotivated 20\% degradation threshold.
  \item Added 4 new citations (Cuoco et al., He \& Garcia, Buda et al., Raghu et al.).
  \item Added rare-class $p$-value footnote to Table~I.
\end{enumerate}

\bigskip
\noindent Sincerely,\\
Jesse Dawson

\end{document}
